\documentclass[12pt]{article}

\usepackage{geometry}
\geometry{margin=1in}
\usepackage{setspace}
%\usepackage{xcolor} % needed for custom colors
\usepackage{hyperref}

\hypersetup{
    colorlinks=true,        % false: boxed links; true: colored links
    linkcolor=blue,         % color of internal links
    citecolor=blue,         % color of citations
    filecolor=magenta,      % color of file links
    urlcolor=blue           % color of external links
}
\usepackage{enumitem}
\usepackage[table]{xcolor}
\usepackage{array}


\usepackage{longtable}
\usepackage{array}
\usepackage[table]{xcolor}
\usepackage{graphicx}
\pagenumbering{gobble}


\definecolor{MLBridge}{HTML}{E40303}       % red
\definecolor{HighDim}{HTML}{FF8C00}        % orange
\definecolor{Nonlinear}{HTML}{FFED00}      % yellow
\definecolor{NeuralNets}{HTML}{008026}     % green
\definecolor{Ensemble}{HTML}{004DFF}       % blue
\definecolor{Heterogeneity}{HTML}{750787}  % violet


\renewcommand{\arraystretch}{1.4}

\usepackage{ragged2e}
\usepackage{longtable}

\setstretch{1.2}

\title{DS-4021: Machine Learning \\ \vspace{0.3cm} \large Syllabus}
\author{School of Data Science, University of Virginia}
\date{}

\begin{document}
\maketitle

\noindent
\textbf{Subject Area and Catalog Number:} Data Science, DS-4021 \\[4pt]
\textbf{Year, Term:} 2026, Fall \\[4pt]
\textbf{Level:} Undergraduate \\[4pt]
\textbf{Credit Type:} Grade (A--F) \\[4pt]
\textbf{Meeting Days:} Tuesday and Thursdays \\[4pt]
\textbf{Meeting Time:} 9:30--10:45 (Section 1) and 11:00--12:15 (Section 2) \\[4pt]
\textbf{Place:} Data Science Building, Room TBD \\[6pt]

\noindent
\textbf{Instructor:} Javier Rasero (\href{mailto:javier.rasero@uv.es}{javier.rasero@uv.es}) \\[2pt]
Office Hours: By appointment \\[2pt]
Location: SDS 335 or Zoom (\href{https://virginia.zoom.us/j/4379931371}{https://virginia.zoom.us/j/4379931371}) \\[8pt]

\noindent
\textbf{Teaching Assistants:} \\[2pt]
Section 1: TBD (\href{mailto:alguien@virginia.edu}{alguien@virginia.edu}) \\[2pt]
Section 2: TBD (\href{mailto:otroalguien@virginia.edu}{otroalguien@virginia.edu}) 

\section*{About the course}

This course continues where Foundations of Machine Learning (DS-3021) left off, moving into more advanced and powerful approaches, particularly neural networks, kernel-based algorithms, and ensemble learning. You will gain a deeper understanding of core ML concepts while comparing methods and discussing how performance, complexity, and predictive power evolve as models become more sophisticated.

The course emphasizes understanding over blind implementation. More than often, you will implement models from scratch rather than relying solely on libraries. You will be challenged to evaluate their assumptions, strengths, and weaknesses. Self-learning and group work are central components, developed through a mix of labs, team-based projects, combining theory, application, discussion and reflection. The course is structured to support students from diverse backgrounds, with ample opportunities for collaboration, TA support, and office hours.

Ultimately, Analytics II seeks to provide a first experience as a real-life practicing data scientist, where advanced models and critical thinking are typically required. It also aims to expand your toolkit with new methods and encourages you to share discoveries with your peers as you improve your ability to generate meaningful insights from data.

\section*{What you will learn}

These are the main learning objectives:

\begin{itemize}[leftmargin=1.5em]
\item Deepen their understanding of machine learning by exploring more advanced and state-of-the-art methods, including support vector machines, ensemble techniques, and neural networks. 
\item Understand the importance of structured workflows in solving data science problems.
\item Gain first exposure to implementing machine learning models from scratch using PyTorch and applying them to real-world datasets.
\item Apply machine learning models real-world datasets using the gold-standard library for machine learning, (scikit-learn.
\item Critically evaluate and compare the performance of different models, considering trade-offs between accuracy, interpretability, and computational cost.
\item Develop and document reproducible data science workflows, including code annotation, version control (GitHub), and final deliverables.
\item Engage in collaborative, project-based learning by working effectively within a group to complete lab assignments and a final project.
\item Demonstrate self-directed learning and adaptability by tackling tasks beyond the core material and seeking solutions independently.
\item Communicate technical findings clearly and effectively, both in written reports and oral presentations.
\item Reflect on individual contributions, strengths, and areas for improvement, incorporating peer and self-critique into a growth-oriented learning process.
\end{itemize}

\section*{How you will learn}

Each week, you will be expected to review short video lectures and complete assigned readings before coming to class. The videos serve as a brief introduction to the main concepts, which will be expanded upon during the first class of the week through a hands-on Jupyter Notebook walkthrough.

Since Data Science is inherently collaborative, nearly every week will include a group-based lab session. These labs are designed not only to help you practice what you have learned that week, but also to push you out of your comfort zone by encouraging you to learn something new. Machine Learning requires a strong element of self-directed learning, and these sessions will help you build that skill. Both the TA and I will be available during and after class to provide any support you need.

Groups will be formed at the beginning of the semester by balancing skill levels, and you will work with the same group throughout the course.

\section*{Evaluation Criteria}

\begin{itemize}

\item \textbf{Labs (50\%):} Throughout the course, and in most weeks, we will have eight in-class graded labs. These are designed to help you either practice the skills presented in class or explore new material related to the week's theme. These are \textbf{group assignments}, and you are expected to work closely with your group mates to divide and complete the different tasks. The main deliverable will be a completed Jupyter Notebook, submitted by the group, which will determine the grade for the assignment. Additionally, each member of the team will submit an individual report describing the specific tasks they contributed to. Each member must submit this individual report separately in order to receive a grade. The deadline for submission of all required materials is typically \textbf{one week} after their publication on Canvas. 

\item \textbf{Final Project (30\%):} The course will culminate in a final project involving work with a dataset of your choice. This is a combined \textbf{group} and \textbf{individual} assignment, and you will work with the same group with which you submitted the lab assignments. As a group, you will submit a final GitHub repository, including all well-annotated code and a final report summarizing your findings. Individually, each member of the team will also submit a GitHub repository containing the code they contributed to different parts of the project, as well as an individual report describing the specific tasks they contributed to. Each member must submit their individual GitHub repository separately in order to receive a grade. \textbf{The final project grade will be the average of the group component and the individual component}. This is an open-ended project designed to allow each group to explore a topic of interest from the semester in greater depth. 

\item \textbf{Final Exam (20\%):} This will be an \textbf{individual} written exam covering material from throughout the course. The exam will take place in class on the date stipulated by the University for this course.

\end{itemize}

\section*{Materials}

\begin{enumerate}[label=\Alph*.]
\item  \href{https://sebastianraschka.com/blog/2022/ml-pytorch-book.html}{Python Machine Learning with PyTorch and Scikit-Learn, 4th Edition}. Freely available on O'Reilly with your UVA account. 
\item \href{https://www.statlearning.com/}{Introduction to statistical learning}. Free PDF. 
\item \href{https://hastie.su.domains/ElemStatLearn/}{The Elements of Statistical Learning: Data Mining, Inference, and Prediction}. Free PDF. 
\item \href{https://www.microsoft.com/en-us/research/people/cmbishop/prml-book/}{Pattern Recognition and Machine Learning}. Free PDF.
\end{enumerate}
\section*{Tech Stack}

\begin{itemize}
\item JupyterLab/Jupyter-notebook 
\item UVA Open OnDemand 
\item Canvas
\end{itemize}

\section*{A few Policies that will Govern the Class}

\textbf{Grading Policies:} Courses carrying a Data Science subject area use the following grading system: A, A-; B+, B, B-; C+, C, C-; D+, D, D-; F.  The symbol W is used when a student officially drops a course before its completion or if the student withdraws from an academic program of the University. \\[2pt]

\noindent
\textbf{Grading Scale:}
\begin{itemize}[leftmargin=1.5em]
\item 93-100 A
\item 90-92 A-
\item 87-89 B+
\item 83-86 B
\item 80-82 B-
\item 77-79 C+
\item 73-76 C
\item 70-72 C-
\item $<$ 70 F
\end{itemize}

\noindent
\textbf{Collaboration and Conflict Resolution:} Should any conflict arise within your group—for example, uneven workload distribution or communication breakdowns—please first make an effort to address the issue respectfully among your group members. If the problem persists or becomes difficult to resolve, group members should reach out to the instructor or TA as early as possible. Addressing concerns early allows us to work together toward a constructive solution before the issue impacts your work or group dynamic. \\[2pt]

\noindent
\textbf{Late work:} Late submissions will NOT be accepted except in cases of emergency, such as illness or family matters. If you anticipate a conflict in advance—for example, due to athletics or another university activity—please reach out ahead of time so we can arrange a suitable alternative for submission. \\[2pt]

\noindent
\textbf{Use of Large Language Model tools:} The use of Large Language Models (e.g., ChatGPT) is acceptable and even encouraged for the coding aspects of the lab assignments and final project, although with caveats. You need to remember that they may make mistakes, so any code generated by this kind of tools should be carefully checked. On the other hand, using LLM in the writing parts of the assignments and final project is discouraged, as it may hinder your own learning. Finally, their use is strictly prohibited during and for quizzes. \\[2pt]

\noindent
\textbf{Religious Accomodations:} It is the University's long-standing policy and practice to reasonably accommodate students so that they do not experience adverse academic consequences when sincerely held religious beliefs or observances conflict with academic requirements. Students who wish to request academic accommodations for a religious observance should submit their request to the instructor by email as far in advance as possible.  If you have questions or concerns about your request, you can contact the University's Office for Equal Opportunity and Civil Rights (EOCR) at \href{mailto:UVAEOCR@virginia.edu}{UVAEOCR@virginia.edu} or (434) 924-3200. Accommodations do not relieve you of the responsibility for the completion of any part of the coursework you miss as the result of a religious observance. \\[2pt]

\noindent
\textbf{University of Virginia Honor System:} By enrolling in this course, all students are expected to uphold the University of Virginia Honor System (\url{https://honor.virginia.edu/course/honor-code}). This includes abiding by the Honor Code in all aspects of the course. \\[2pt]

\noindent
\textbf{Special Needs:}  The University of Virginia accommodates students with disabilities. Any SCPS student with a disability who needs accommodation (e.g., in arrangements for seating, extended time for examinations, or note-taking, etc.), should contact the Student Disability Access Center (SDAC) and provide them with appropriate medical or psychological documentation of his/her condition. Once accommodations are approved, just follow up with me concerning any logistics and implementation of accommodations.  Please try to make accommodations for test-taking at least 14 business days in advance of the date of the test(s). Students with disabilities are encouraged to contact the SDAC: 434-243-5180/Voice, 434-465-6579/Video Phone, 434-243-5188/Fax. Further policies and statements are available at \url{www.virginia.edu/studenthealth/sdac/sdac.html}.

\newpage

\section*{Class Schedule (subject to change)}
\renewcommand{\arraystretch}{1.25}
\begin{longtable}{|>{\raggedright}p{2.5cm}|>{\raggedright}p{5cm}|>{\raggedright}p{5cm}|>{\raggedright\arraybackslash}p{3cm}|}
\hline
\textbf{Week} & \textbf{Tuesday Class} & \textbf{Thursday Class} & \textbf{Suggested Readings} \\
\noalign{\global\arrayrulewidth=1.2pt}
\hline
\noalign{\global\arrayrulewidth=0.4pt}
\multicolumn{4}{|c|}{\cellcolor{MLBridge!70}\textbf{Module 1: ML1 Bridge}} \\
\noalign{\global\arrayrulewidth=1.2pt}
\hline
\noalign{\global\arrayrulewidth=0.4pt}
\rowcolor{MLBridge!25}
25-27 August & Course Overview. \newline Skill Assessment (for Group Formation) & Feature wrangling: transformations, missingness, scaling, selection. & \\
\hline
\rowcolor{MLBridge!25}
1-3 September & Mastering model evaluation and optimization  & Introduction to PyTorch & \\
\noalign{\global\arrayrulewidth=1.2pt}
\hline
\noalign{\global\arrayrulewidth=0.4pt}
\multicolumn{4}{|c|}{\cellcolor{HighDim!70}\textbf{Module 2: Dealing with High Dimensional Data}} \\
\noalign{\global\arrayrulewidth=1.2pt}
\hline
\noalign{\global\arrayrulewidth=0.4pt}
\rowcolor{HighDim!25}
8-10 September &    Regularized Regression through PyTorch and Stochastic Gradient Descent &  Principal Component Analysis as a preprocessing step. \newline Introduction to Pipelines &
\\
\hline
\rowcolor{HighDim!25}
15-17 September & {\bf Lab}: Partial Least Squares Regression (PLS) & \cellcolor{Nonlinear!25} Support Vector Machine I: The Linear Case & \\
\noalign{\global\arrayrulewidth=1.2pt}
\hline
\noalign{\global\arrayrulewidth=0.4pt}
\rowcolor{Nonlinear}
\multicolumn{4}{|c|}{\cellcolor{Nonlinear!70}\textbf{Module 3: Exploiting Non-Linear Relationships in Data}} \\
\noalign{\global\arrayrulewidth=1.2pt}
\hline
\noalign{\global\arrayrulewidth=0.4pt}
\rowcolor{Nonlinear!25}
22-24 September & Support Vector Machine II: The Non-Linear case through the Kernel Trick & {\bf Lab}: Predicting Seismic Building Response with Support Vector Regrssion & A. Chapter 3: {\small Maximum margin classification with support vector machines}\\
\hline
\rowcolor{Nonlinear!25}
29 September-1 Octobe &  Kernels are EVERYWHERE & {\bf Lab}: Non-linear dimensionality reductions: t-SNE, and UMAP & A. Chapter 3: {\small Solving nonlinear problems using a kernel SVM}\\
\noalign{\global\arrayrulewidth=1.2pt}
\hline
\noalign{\global\arrayrulewidth=0.4pt}
%\rowcolor{NeuralNets!45}
\multicolumn{4}{|c|}{\cellcolor{NeuralNets!70}\textbf{Module 4: Neural Networks}} \\
\noalign{\global\arrayrulewidth=1.2pt}
\hline
\noalign{\global\arrayrulewidth=0.4pt}
\rowcolor{NeuralNets!25}
6-8 October & Reading Day (No Class) & Neural Networks I: Introduction \& Multi-Layer Perceptrons (MLPs) & \\
\hline
\rowcolor{NeuralNets!25}
13-15 October &  {\bf Lab}: Manual implementation of MLPs for Classification and Regression using toy data & Neural Networks II: Deeper into model building, training and opmitization. & A. Chapter 11: {\small  Modeling complex functions with artificial neural networks}\\
\hline
\rowcolor{NeuralNets!25}
20-22 October & {\bf Lab}: Application of  Neural Networks to Real Data & Neural Networks III: Improving generalization & A. Chapter 13: {\small Making model building more flexible with nn.Module} \\
\hline
\rowcolor{NeuralNets!25}
27-29 October & {\bf Lab}: Implementing an autoencoder for non-linear dimensionality reduction &  \cellcolor{Ensemble!25} Bagging: Bootstrap Aggregating and Random Forest  & \\
\noalign{\global\arrayrulewidth=1.2pt}
\hline
\noalign{\global\arrayrulewidth=0.4pt}
\multicolumn{4}{|c|}{\cellcolor{Ensemble!70}\textbf{Module 5: Ensemble Methods \& Meta-Learning}} \\
\noalign{\global\arrayrulewidth=1.2pt}
\hline
\noalign{\global\arrayrulewidth=0.4pt}
\rowcolor{Ensemble!25}
3-5 November & Election day (No class) & Boosting: Adaptive Boosting and Gradient Boosting  & \\
\hline
\rowcolor{Ensemble!25}
10-12 November &  {\bf Lab}: Classifying Internet of Things Device Types with Bagging and Boosting  & \cellcolor{Heterogeneity!55} Clustering I: k-Means and Gaussian Mixture Models &\\
\noalign{\global\arrayrulewidth=1.2pt}
\hline
\noalign{\global\arrayrulewidth=0.4pt}
\multicolumn{4}{|c|}{\cellcolor{Heterogeneity!70}\textbf{Module 6: Tackling Population Heterogeneity}} \\
\noalign{\global\arrayrulewidth=1.2pt}
\hline
\noalign{\global\arrayrulewidth=0.4pt}
\rowcolor{Heterogeneity!25}
17-19 November &  Clustering II: Hierarchical Clustering and DBSCAN & Clustering III: Validation & \\
\hline
\rowcolor{Heterogeneity!25}
26-28 November &   {\bf Lab}: clustering &  Thanksgiving (No class) & \\
\hline
1-3 December & \multicolumn{3}{|c|}{Final Project Time} \\
\hline
\end{longtable}
\end{document}
